% =====================================================================
% Beyond the Human Eye: A Three-Layer Verification Architecture for
% Large-Scale Scholarly Wikidata Item Creation
%
% Anonymized draft for ESWC/ISWC (double-blind) — llncs single-column.
% NOTE TO AUTHORS (remove before submission):
%  - \cite entries marked VERIFY-ARXIV need bibliographic confirmation.
%  - Unblinded copy (repo URLs, institution) kept out of version control.
%  - Countersign subset (200 rows) pending curator labels; Section 7
%    currently reports the Claude-labeled gold with this disclosed.
% =====================================================================
\documentclass{llncs}
\usepackage[T1]{fontenc}
\usepackage[utf8]{inputenc}
\usepackage{amsmath}
\usepackage{booktabs}
\usepackage{graphicx}
\usepackage{xcolor}
\usepackage{url}
\usepackage{hyperref}

\title{Beyond the Human Eye: A Three-Layer Verification Architecture for
Large-Scale Scholarly Wikidata Item Creation}

\author{Anonymous}
\institute{Anonymous affiliation}

\begin{document}
\maketitle

\begin{abstract}
Publishing scholarly catalogue records as Wikidata and project-Wikibase
items makes identity-bearing claims public at a scale that outgrows human
review. We present the verification architecture of a production pipeline
that converts Hebrew-manuscript MARC records into an ontological RDF graph
and then into Wikibase items, and we ask which machine judge may act in
place of the curator's eye. We compare three layers on identical fixtures:
a deterministic rule engine (30 checks, including fail-closed live API
rules), a generative LLM judge driven by incident-evolved rubrics, and a
decision model trained with Reinforcement Learning from Contrastive
Distillation (RLCD) that returns calibrated, typed answers instead of
generated text. We contribute (i) the three-layer architecture with
responsibilities assigned by measurement; (ii) a certification
methodology---gold construction, an \emph{unsafe-approve} metric, accuracy
and determinism gates, and a tuning loop that converts every misjudgement
into either a deterministic check or a rubric-encoded rule; and (iii) an
empirical evaluation on 1{,}238 gold-labeled items across six verification
channels. Our results show that (a) deterministic code, not model judgment,
should own mechanical artifact detection (moving four artifact classes from
the model to code removed 10 of 11 unsafe approvals); (b) LLM judges are
inconsistent on identical defects, while atomic decomposition of rubric
axes into typed questions makes the RLCD judge uniform; (c) calibrated
confidence gating converts model uncertainty into curator review rather
than silent failures; and (d) an escalation policy---RLCD judge primary,
LLM as exception handler---reduces LLM workload to roughly 15\% of rows
while wall-clock verification drops from hours to under two minutes per
500 items. Neither judge yet passes the zero-unsafe certification gate
(0.65\% unsafe-approve over the full gold for the RLCD judge, 2.4\% on the
LLM judge's overlapping subset), so the system keeps human approval as the
final authority and assigns the LLM to confirming rows the decision model
is not confident about.
\end{abstract}

\section{Introduction}

Cultural-heritage institutions increasingly publish their catalogues as
linked data. When the target is Wikidata or a project Wikibase, item
creation is the risk surface: items are public, they assert identity
(``this work is by this person''), and errors are costly to detect and
undo after third parties consume them. In our domain---Hebrew manuscripts
catalogued in MARC, projected through a CIDOC-CRM-based ontology into a
project Wikibase and onward to public Wikidata under the WikiProject
Manuscripts data model---a single run of the pipeline can propose thousands
of items. Reviewing each item with the human eye does not scale.

This paper reports the verification architecture that replaced per-item
human review in production, and the methodology we use to decide when a
machine judge may act. The architecture has three layers:

\begin{enumerate}
\item a \textbf{deterministic rule engine} (30 rules: SHACL wrappers,
  export-quality codes, MARC grounding, duplicate probes, fail-closed live
  API checks);
\item a \textbf{generative LLM judge} (Gemini with an agentic tool loop;
  Qubrid-hosted Kimi K2.5 and DeepSeek in linear mode) driven by rubrics
  that encode production incidents as decision procedures; and
\item a \textbf{RLCD decision model} (TypeSafe Jev) that evaluates typed
  questions against the evidence and returns calibrated answers with
  confidence instead of generated text.
\end{enumerate}

The central engineering question is not which model is ``best'' but
\emph{which layer may make which decision, with what evidence}. Our answer
is a certification methodology: a gold set with the unsafe-approve metric,
accuracy and determinism gates, and a tuning loop in which every gold miss
becomes either a deterministic check or a rubric-encoded question rule. We
report measurements from production data (run of 68 manuscripts, 1{,}238
gold-labeled verification rows, full populations for the extraction and
Wikidata channels) and a double-blind-comparable evaluation design.

\paragraph{Contributions.}
\begin{itemize}
\item A three-layer verification architecture for large-scale item
  creation, with per-layer responsibility assignment justified by
  measurement (Sect.~\ref{sec:layers}).
\item A certification methodology for machine judges: gold construction,
  the unsafe-approve metric, acceptance gates, double-run determinism, and
  a miss-to-fix tuning loop (Sect.~\ref{sec:cert}).
\item An empirical study comparing the LLM judge and the RLCD judge on
  identical fixtures against a common gold set, including an ablation
  isolating the value of deterministic gates (Sect.~\ref{sec:eval}).
\item A production escalation policy: decision model primary, LLM as
  exception handler, no automatic approval before certification
  (Sect.~\ref{sec:policy}).
\end{itemize}

\section{Background and Related Work}
\label{sec:related}

\paragraph{Wikidata and Wikibase quality.}
Wikidata's collaborative model~\cite{vrandecic2014wikidata} has made item
quality a research topic; quality assessments range from completeness and
consistency~\cite{piscopo2019quality} to schema-based validation. Project
Wikibases inherit the same risks at creation time, before community
checking exists. SHACL~\cite{knublauch2017shacl} provides closed-shape
validation and gates our upload path, but many failure modes of item
creation (misattribution, wrong entity links, generated-text artifacts,
identity contamination) are semantic and out of SHACL's reach.

\paragraph{LLM-as-judge.}
Using LLMs as evaluators is established~\cite{zheng2023judging}, with known
weaknesses: position bias, verbosity bias, and---most relevant here---%
\emph{self-consistency}: identical inputs can receive different
verdicts~\cite{wang2022selfconsistency}. Our measurements confirm the
consistency problem in a verification (not preference) setting: an LLM
judge flagged the same description artifact on some rows and missed it on
twelve identical rows. Self-consistency voting~\cite{wang2022selfconsistency}
is the standard mitigation; we show that \emph{atomic decomposition} of the
rubric into independent typed questions achieves uniformity at lower cost.

\paragraph{Calibration and abstention.}
Whether language models ``know what they don't know'' is studied
extensively~\cite{kadavath2022language}. Our architecture operationalises
this: the RLCD judge returns calibrated probabilities and a confidence
score, and the escalation policy treats low-confidence verdicts as
\emph{review}, never as approval or rejection. This turns abstention---%
long identified as the safe failure mode---into the system's primary
safety mechanism.

\paragraph{RLCD and decision models.}
Reinforcement Learning from Contrastive
Distillation~\cite{yang2023rlcd} trains models to prefer contrastive
chosen/rejected pairs; TypeSafe's System One models apply the paradigm to
produce a non-generative judge that evaluates typed questions
(\texttt{Choice}, \texttt{Score}, binary \texttt{Noul}) against a shared
state and returns probabilities over the supplied options. To our
knowledge this is the first reported use of such a decision model as the
primary quality judge in a semantic-web publication pipeline.

\paragraph{Human-in-the-loop escalation.}
Escalation from automated judgment to human review is a staple of
deployed ML systems. Our contribution is the specific policy: the cheap,
fast, calibrated judge runs on every row; the expensive generative judge
runs only on rows the primary judge is not highly confident about; and
certification gates decide when either judge may auto-approve (today:
never).

\section{System Background}
\label{sec:system}

The pipeline has five curator-reviewed stages: (1)~AI extraction of named
entities and genres from MARC records; (2) authority enrichment (Mazal/NLI,
VIAF, KIMA, Wikidata reconciliation, fail-closed matching); (3) RDF graph
construction over a CIDOC-CRM-based Hebrew-manuscripts ontology; (4)
projection to a project Wikibase (HMO Wikibase Studio); and (5) guarded
writes to public Wikidata following the WikiProject Manuscripts data
model. Every curator mutation is event-versioned, and all slow work runs
as heartbeated background jobs with a Postgres read-model behind a
Redis-first cache stack.

Item creation happens at stages 4 and 5. Each proposed item carries a
compact evidence pack: a filtered MARC slice keyed by control numbers,
authority packs (VIAF, Mazal), existing-Wikidata and duplicate-probe
results, per-claim provenance, and validator output. Verdicts are cached
by content fingerprint, so re-verification after edits is incremental.

\section{Layer 1: The Deterministic Rule Engine}
\label{sec:rules}

The rule engine runs every deterministic check the AI judges also see,
without AI. Its 30 rules fall into three families:

\begin{enumerate}
\item \textbf{Wrapped validators}: SHACL violation/error severities
  (upload-gating) and one rule per export-quality issue code (13 codes,
  e.g.\ blank-node exports, malformed labels, generic fallback
  descriptions), so each code is individually blockable by a curator.
\item \textbf{Deterministic checks}: class/source presence, claim datatype
  shapes, MARC grounding (the item's identity control number must appear
  in its own local identifier; descriptions must cite the item's own
  manuscript), within-run duplicate labels, language hygiene.
\item \textbf{Fail-closed API rules}: live Wikibase liveness and label
  drift, Wikidata QID liveness, and a duplicate-probe search. Without API
  access they return \texttt{error}---never \texttt{pass}.
\end{enumerate}

Results follow an explicit contract (\texttt{pass}/\texttt{fail}/%
\texttt{not\_relevant}/\texttt{error}); rule results are advisory except
SHACL, which gates upload. Crucially for this paper, the rule engine owns
the \emph{mechanical artifact} class. During certification (Sect.~%
\ref{sec:eval:ablation}) we found that encoding four mechanical checks in
code---per-text parenthesis balance, identity-control-number vs.\
description mismatch, quote-collision garbles in Hebrew abbreviations, and
truncated folio ranges---removed 10 of 11 unsafe approvals that model
judgment alone produced. Deterministic code is also where the LLM judge
itself is weakest: the same generated-template defect was flagged on some
rows and missed on twelve identical rows by the LLM, while the rule fired
uniformly.

\section{Layer 2: The LLM Judge}
\label{sec:llm}

The vendored judge agent reads the evidence pack and returns a JSON
verdict with three axes (\texttt{name\_ok}, \texttt{type\_ok},
\texttt{role\_ok}) and an \texttt{overall} of \texttt{full}/%
\texttt{partial}/\texttt{fail}/\texttt{abstain}. Rubrics are
incident-evolved documents: each production incident (hallucinated
institutions, invented ``missing units'', cross-record contamination,
wrong public-QID links) closed with a rubric rule plus a regression
example. The judge runs in two modes: an agentic tool loop (Gemini) that
can pull full MARC records and authority packs on demand, and a linear
single-shot mode for hosted models (Qubrid Kimi K2.5, DeepSeek V4).

Two structural costs dominate. First, \emph{consistency}: the judge
re-derives its own search of the evidence at call time, so identical
defects receive different verdicts across rows. Second, \emph{latency and
cost}: the judge generates a full reasoning-plus-JSON response per row
(9.2--22.2\,s per item in our measurements; 313k input tokens for 40 items,
14\,m\,48\,s wall). A 500-item verification pass takes roughly three hours.

\section{Layer 3: The RLCD Decision Model}
\label{sec:jev}

The third layer reformulates judging as \emph{evaluation against typed
questions}. The judge prompt (identical to the LLM judge's, including
rubric, prediction block, MARC context, grounding signal, and skill
packs) becomes the \emph{state}; each rubric axis becomes a typed
question; and the model returns calibrated answers instead of generated
text:

\begin{itemize}
\item \textbf{Choice} per axis: \texttt{name\_ok}, \texttt{type\_ok},
  \texttt{role\_ok} each become a Choice question over
  \{yes, partial, no\} (plus \texttt{not\_applicable} for structural
  entities), with the rubric's per-entity-type exceptions encoded in the
  question criteria.
\item \textbf{Binary Noul per claim}: each present statement (P31, P1476,
  P2093, P50, P8189, \dots) is one question---``is this statement
  supported by at least one evidence channel?''---answered with a
  probability. Questions are evaluated in parallel and independently;
  adding checks does not lengthen the response.
\item \textbf{Atomic quality questions}: label/description quality
  (system-label acceptance per entity type) and a generation-artifact
  check read ``character by character.''
\end{itemize}

Three design consequences follow. First, \texttt{overall} is \emph{never}
asked of the model: it is computed in code from the check states, so the
rubric's universal table (and its fail-over-partial tiebreaker) is
deterministic. Second, \emph{confidence gating} becomes the safety
mechanism: a blocking \texttt{no} below a calibrated threshold is routed
to review rather than issued as a hard fail (Sect.~\ref{sec:eval}).
Third, every row yields a structured \emph{check list}
(rule id, state, message, evidence, blocking flag) rendered in the same
visual language as the rule-engine panel, so curators see exactly which
checks failed or partially failed. Deterministic code gates (validator
errors, mechanical artifacts) are applied after the model answers and can
never be softened by confidence.

\section{Certifying the Judges}
\label{sec:cert}

No judge is trusted by assertion. Each judge must pass a certification
before it may act on a channel:

\paragraph{Gold set.}
1{,}238 rows labeled against the same evidence the judges see: full
populations for the four extraction evaluators (507 rows) and the Wikidata
channel (232 rows), plus a 500-row stratified sample of the 1{,}949-item
Wikibase projection (exportable items only; non-exportable drafts
excluded). Labels cover \texttt{overall} plus the three axes, carry a
labeler confidence and an escalation flag; 14 rows are flagged for curator
review. A 200-row stratified subset (67/66/67 full/partial/fail) is
prepared for curator countersigning; the results in this paper use the
pre-countersign labels and this is disclosed as a limitation.

\paragraph{Unsafe-approve metric.}
The only error direction that can ship a bad item is a \texttt{full}
verdict on a gold-\texttt{partial}/\texttt{fail} row. We therefore report
\emph{unsafe-approve} separately from accuracy, and its gate is zero.
\emph{Safe-miss} (a \texttt{fail} verdict on a gold-\texttt{full} row)
costs review effort but cannot ship an error; we report it to separate
precision loss from safety loss.

\paragraph{Gates.}
Certification requires, per channel: accuracy $\geq 0.98$ against gold,
unsafe-approve $= 0$, and double-run verdict stability $\geq 0.99$.
Below the gates, judges may triage but not approve.

\paragraph{Tuning loop.}
Every certification miss is converted into a permanent fix: mechanical
artifacts become deterministic checks; domain nuances (dictation notes vs.\
copyists; father/teacher vs.\ owner in provenance notes; boundary-noise
spans; claim references establishing provenance) become encoded question
criteria; and error-severity validator findings become code-level gates
that confidence can never soften. The loop is the paper's practical
recipe: \emph{each gold miss must end in code or in a rule, never in a
prompt hope}.

\section{Evaluation}
\label{sec:eval}

All experiments use production data (one pipeline run; 68 manuscripts;
1{,}949 projected entities; 232 proposed Wikidata items; 507 extraction
predictions). Judges see byte-identical evidence. The RLCD judge processes
items concurrently ($\sim$5 in flight); the LLM judge runs linearly at
30\,rpm.

\subsection{Head-to-head on identical fixtures (E1)}
\label{sec:eval:e1}

On 40 items per channel judged by both systems on identical fixtures, the
RLCD judge averaged 0.37--0.61\,s per item (a 500-item scope completes in
39--77\,s wall) versus 9.2--22.2\,s per item for the LLM (40 HMO items in
14\,m\,48\,s; a 500-item scope $\approx$ 3\,h\,05\,m): a
\textbf{100--250$\times$ wall-clock advantage} for the RLCD judge.
Overall agreement between the judges was 0.85--0.975 depending on channel;
crucially, \emph{every} RLCD disagreement with the LLM was in the
conservative direction (route to review), and the LLM missed 12
identical template defects that the RLCD judge flagged uniformly via its
atomic artifact question.

\subsection{Certification against gold (E2)}
\label{sec:eval:e2}

Table~\ref{tab:cert} reports the post-tuning certification numbers.

\begin{table}[t]
\centering
\small
\begin{tabular}{lrrrrr}
\toprule
Channel & $n$ & Accuracy & Unsafe-approve & Safe-miss & Cost/item \\
\midrule
Person NER & 204 & 0.809 & \textbf{0} & 9 & \$0.21/1k \\
Contents NER & 150 & 0.940 & 1 & 0 & \$0.24/1k \\
Provenance NER & 60 & 0.817 & 2 & 2 & \$0.16/1k \\
Genre classifier & 93 & 0.914 & 2 & 2 & \$0.15/1k \\
Wikibase items & 500 & 0.856 & 1 & 4 & \$0.34/1k \\
Wikidata items & 231 & 0.671 & 2 & 0 & \$0.63/1k \\
\midrule
Total & 1{,}238 & 0.826 & 8 (0.65\%) & 17 & \$0.42 total \\
\bottomrule
\end{tabular}
\caption{RLCD judge vs.\ the 1{,}238-row gold (full populations for the
extraction and Wikidata channels; stratified 500 for the Wikibase
projection). Accuracy gate is 0.98; the unsafe-approve gate is 0. Costs at
the published rate of \$42 per Btok (\$0.042/Mtok input, output tokens
free; console-verified: 54.2M tokens billed \$2.20), reported per 1{,}000
items. The full 1{,}238-row certification cost \$0.42.}
\label{tab:cert}
\end{table}

For reference, the LLM judge scored on the overlapping gold subset
($n{=}123$) reaches 0.902 accuracy with 3 unsafe approvals (2.4\% on that
subset); its subsets under-represent the hardest classes, so the numbers
are indicative rather than directly comparable. On the wikidata
identifier-less cohort---the channel's dominant failure class---both
judges fully agree with gold (63/63) once the deterministic validator gate
is in place.

\subsection{Ablation: model-only vs.\ model + code gates (E3)}
\label{sec:eval:ablation}

Table~\ref{tab:ablation} isolates the deterministic gates' contribution on
the 500-item Wikibase channel and the Wikidata channel. Moving mechanical
artifacts and validator severities from model judgment to code removed
10 of 11 unsafe approvals and made the dominant failure class
(\texttt{NO\_IDENTIFIER}, error severity) exactly reproducible (63/63).

\begin{table}[t]
\centering
\small
\begin{tabular}{lccc}
\toprule
 & Model-only & + code gates & $\Delta$ \\
\midrule
Wikibase unsafe-approve (/500) & 11 & 1 & $-10$ \\
Wikidata fail-class recall vs.\ gold & 34/63 & 63/63 & $+29$ \\
Person unsafe-approve & 1 & 0 & $-1$ \\
\bottomrule
\end{tabular}
\caption{Ablation of deterministic gates (validator severity in code;
artifact checks in code; rubric nuances as encoded question criteria).}
\label{tab:ablation}
\end{table}

\subsection{Escalation policy (E4)}
\label{sec:eval:e4}

The production policy is: RLCD judge primary; LLM confirms any row that is
not a high-confidence \texttt{full} (measured 15--25\% of rows); neither
judge auto-approves before certification passes. Under this policy the LLM
workload drops to roughly 15\% of today's, the LLM's unsafe approvals are
confined to the subset it confirms, and a 500-item verification completes
in under two minutes instead of three hours. Because verdict caching is
content-addressed, the escalation path reuses the same session, cache, and
persistence machinery as the existing single-judge flow.

\subsection{Determinism}

A double-run of the 500-item Wikibase channel measured 0.966 overall and
0.99 per-axis verdict stability for the RLCD judge; all unstable verdicts
flipped between \texttt{full} and \texttt{partial}---the review direction.
This is below our 0.99 gate at the overall level and motivates
majority-voting borderline axes in production (future work).

\subsection{Cost}

The RLCD judge costs \$0.15--0.63 per 1{,}000 items at the published rate
of \$42 per Btok (\$0.042/Mtok input; output tokens free;
console-verified: 54.2M tokens billed \$2.20). The complete 1{,}238-row
certification cost \$0.42. The LLM judge consumed 7.8k--10.7k input and
238--422 output tokens per item; its cost is
$c_{\text{item}} = t_{\text{in}} \cdot r_{\text{in}} + t_{\text{out}}
\cdot r_{\text{out}}$ with the provider rate as parameter. At any
plausible hosted-LLM rate the generative judge costs more per item than
the RLCD judge---and is 50--250$\times$ slower. With the corrected price,
the RLCD judge dominates on speed, cost, and structure simultaneously;
the generative judge's remaining advantages are its free-text audit trail
and its (subset-measured) accuracy edge.

\section{Production Policy}
\label{sec:policy}

The deployed configuration follows from the measurements:
\begin{enumerate}
\item \textbf{RLCD judge primary} on every verification channel, emitting
  the structured check list per item;
\item \textbf{LLM as exception handler} for rows that are not
  high-confidence passes;
\item \textbf{no automatic approval} by any judge: the certification gate
  (zero unsafe-approve, accuracy $\geq 0.98$, determinism $\geq 0.99$)
  must pass on the full gold first, and SHACL remains the only upload
  gate;
\item curators act on the union of rule-engine failures and AI check
  lists, with identical UI semantics across layers.
\end{enumerate}

\section{Limitations}
\label{sec:limits}

The gold set was labeled by a strong LLM against the same evidence the
judges see; curator countersigning is prepared (200 stratified rows) but
pending, and one artifact class showed labeler disagreement that caps
measurable accuracy ($\approx$0.90 reconciled on the Wikibase channel).
The domain is a single manuscript pipeline; replication on other
catalogues is future work. The RLCD model is English-primary and
Hebrew-heavy states are not independently certified at scale. Overall
determinism (0.966) is below the 0.99 gate. Finally, the LLM judge's
full-gold certification was not run (a multi-hour job); its reported
numbers are subset-based.

\section{Conclusion and Future Work}

Large-scale scholarly item creation for Wikidata and project Wikibases can
shed the per-item human eye---but only if judge authority is earned by
certification, mechanical truth is pushed into code, semantic judgment is
decomposed into atomic calibrated questions, and uncertainty is routed to
humans rather than resolved silently. Our measurements show the RLCD
decision model is the right primary judge on speed, structure, and
safety-direction, with the LLM judge as a confirming exception handler;
neither judge is yet certifiable for autonomous approval, and the
certification methodology itself---gold, gates, and the miss-to-fix
loop---is the paper's reusable contribution. Future work includes
curator-countersigned certification, self-consistency voting for
borderline axes, multi-domain replication, and learned question policies
that allocate model questions only where code cannot decide.

\paragraph{Data availability.}
All measurements derive from the production runs described in
Sect.~\ref{sec:system}; anonymized gold labels, per-run reports, and the
harness scripts are available from the authors.

\begin{thebibliography}{9}

\bibitem{vrandecic2014wikidata}
D. Vrande\v{c}i\'c, M. Kr\"otzsch.
Wikidata: A Free Collaborative Knowledgebase.
\emph{Communications of the ACM} 57(10):78--85, 2014.

\bibitem{piscopo2019quality}
A. Piscopo, E. Simperl.
What We Talk About When We Talk About Wikidata: Quality and Use.
\emph{Proc.\ WWW}, 2019.

\bibitem{knublauch2017shacl}
H. Knublauch, D. Kontokostas.
Shapes Constraint Language (SHACL).
W3C Recommendation, 2017.

\bibitem{zheng2023judging}
L. Zheng et al.
Judging LLM-as-a-Judge with MT-Bench and Chatbot Arena.
\emph{Proc.\ NeurIPS Datasets and Benchmarks}, 2023.

\bibitem{wang2022selfconsistency}
X. Wang et al.
Self-Consistency Improves Chain of Thought Reasoning in Language Models.
arXiv:2203.07181, 2022.

\bibitem{kadavath2022language}
S. Kadavath et al.
Language Models (Mostly) Know What They Know.
arXiv:2207.05221, 2022.

\bibitem{yang2023rlcd}
% VERIFY-ARXIV-ID before submission
Y. Yang et al.
RLCD: Reinforcement Learning from Contrastive Distillation for Language
Model Alignment.
arXiv preprint, 2023.

\end{thebibliography}

\end{document}
